% GENERATED FILE. Edit canonical lessons, then run npm run generate:books.
% canonical-source-hash: fnv1a64:2b1c277795bc53f8
% canonical-lessons: FA-C450-markaz, FA-C450-hume, FA-C450-borj, FA-C450-mojassame, FA-C450-favvare

\chapter{The town centre}
\label{ch:fa-the-town-centre}
\hlchaptermodality{450}

\begin{chapteropening}
\textbf{By the end of this chapter:} \emph{I can describe the town centre: the centre, the outskirts, a tower, a statue and a fountain.}

Complete the last lesson of chapter 450: I can describe the town centre: the centre, the outskirts, a tower, a statue and a fountain.
\end{chapteropening}

\section[\texorpdfstring{\fa{مرکز} (markaz)}{مرکز (markaz)}]{\fa{مرکز} (markaz) — a centre}

\label{lesson:FA-C450-markaz}



\begin{quote}
Before the new one: say the Persian for to pass, then the Persian for the back streets.
\end{quote}

\subsection*{You'll want to know: \fa{مرکز}}
\textbf{\fa{مرکز}} (\emph{markaz}) — \textquotedblleft{}a centre\textquotedblright{}.

\begin{quote}
\fa{هتل} \fa{ما} \fa{در} \fa{مرکز} \fa{شهر} \fa{است}.
\end{quote}

\begin{quote}
\emph{Hotel-e mâ dar markaz-e shahr ast.}
\end{quote}

\begin{quote}
Our hotel is in the town centre.
\end{quote}

Arabic plural \emph{marâkez}. \emph{Markaz-e kharid} — a shopping centre.

\begin{cousinweb}
\emph{markaz} was borrowed from Arabic \emph{markaz} 'centre', which EtymDB lists.
\end{cousinweb}

\begin{grammarlens}[title={What you've built}]
One more word for telling, explaining and describing at B1.
\end{grammarlens}

\subsection*{Your turn}
\begin{itemize}
\raggedright
  \item \emph{Say it:} \emph{markaz}
  \item \emph{Say it:} \emph{markaz}, once more
  \item \emph{Say it:} a sentence of your own with \emph{markaz}
  \item \emph{Recall:} say the Persian for to pass, then the Persian for the back streets, then say \emph{markaz} again
\end{itemize}

\begin{tcolorbox}[breakable,colback=teal!4,colframe=teal!35!black,title={Before you move on}]
What does \fa{مرکز} mean? (\textquotedblleft{}A centre\textquotedblright{}.) Say it once more.
\end{tcolorbox}
\section[\texorpdfstring{\fa{حومه} (hume)}{حومه (hume)}]{\fa{حومه} (hume) — the outskirts, the suburbs}

\label{lesson:FA-C450-hume}



\begin{quote}
Before the new one: say the Persian for the back streets, then the Persian for a centre.
\end{quote}

\subsection*{You'll want to know: \fa{حومه}}
\textbf{\fa{حومه}} (\emph{hume}) — \textquotedblleft{}the outskirts, the suburbs\textquotedblright{}.

\begin{quote}
\fa{خانواده} \fa{ما} \fa{در} \fa{حومه} \fa{شهر} \fa{زندگی} \fa{میکند}.
\end{quote}

\begin{quote}
\emph{Khânevâde-ye mâ dar hume-ye shahr zendegi mikonad.}
\end{quote}

\begin{quote}
Our family lives on the outskirts of the town.
\end{quote}

Usually \emph{hume-ye shahr}.

\begin{cousinweb}
\emph{hume} was borrowed from Arabic \emph{ḥawma}, which EtymDB lists; it is usually glossed 'the area around, the edge'.
\end{cousinweb}

\begin{grammarlens}[title={What you've built}]
One more word for telling, explaining and describing at B1.
\end{grammarlens}

\subsection*{Your turn}
\begin{itemize}
\raggedright
  \item \emph{Say it:} \emph{hume}
  \item \emph{Say it:} \emph{hume}, once more
  \item \emph{Say it:} a sentence of your own with \emph{hume}
  \item \emph{Recall:} say the Persian for the back streets, then the Persian for a centre, then say \emph{hume} again
\end{itemize}

\begin{tcolorbox}[breakable,colback=teal!4,colframe=teal!35!black,title={Before you move on}]
What does \fa{حومه} mean? (\textquotedblleft{}The outskirts, the suburbs\textquotedblright{}.) Say it once more.
\end{tcolorbox}
\section[\texorpdfstring{\fa{برج} (borj)}{برج (borj)}]{\fa{برج} (borj) — a tower; a tower block}

\label{lesson:FA-C450-borj}



\begin{quote}
Before the new one: say the Persian for a centre, then the Persian for the outskirts.
\end{quote}

\subsection*{You'll want to know: \fa{برج}}
\textbf{\fa{برج}} (\emph{borj}) — \textquotedblleft{}a tower; a tower block\textquotedblright{}.

\begin{quote}
\fa{از} \fa{بالای} \fa{برج} \fa{تمام} \fa{شهر} \fa{را} \fa{دیدیم}.
\end{quote}

\begin{quote}
\emph{Az bâlâ-ye borj tamâm-e shahr râ didim.}
\end{quote}

\begin{quote}
From the top of the tower we saw the whole town.
\end{quote}

Arabic plural \emph{boruj}. Also a sign of the zodiac.

\begin{cousinweb}
\emph{borj} was borrowed from Arabic \emph{burj} 'tower', which EtymDB traces through Syriac to Greek \emph{pýrgos} 'tower'.
\end{cousinweb}

\begin{grammarlens}[title={What you've built}]
One more word for telling, explaining and describing at B1.
\end{grammarlens}

\subsection*{Your turn}
\begin{itemize}
\raggedright
  \item \emph{Say it:} \emph{borj}
  \item \emph{Say it:} \emph{borj}, once more
  \item \emph{Say it:} a sentence of your own with \emph{borj}
  \item \emph{Recall:} say the Persian for a centre, then the Persian for the outskirts, then say \emph{borj} again
\end{itemize}

\begin{tcolorbox}[breakable,colback=teal!4,colframe=teal!35!black,title={Before you move on}]
What does \fa{برج} mean? (\textquotedblleft{}A tower; a tower block\textquotedblright{}.) Say it once more.
\end{tcolorbox}
\section[\texorpdfstring{\fa{مجسمه} (mojassame)}{مجسمه (mojassame)}]{\fa{مجسمه} (mojassame) — a statue}

\label{lesson:FA-C450-mojassame}



\begin{quote}
Before the new one: say the Persian for the outskirts, then the Persian for a tower.
\end{quote}

\subsection*{You'll want to know: \fa{مجسمه}}
\textbf{\fa{مجسمه}} (\emph{mojassame}) — \textquotedblleft{}a statue\textquotedblright{}.

\begin{quote}
\fa{وسط} \fa{میدان} \fa{یک} \fa{مجسمه} \fa{بزرگ} \fa{است}.
\end{quote}

\begin{quote}
\emph{Vasat-e meydân yek mojassame-ye bozorg ast.}
\end{quote}

\begin{quote}
There is a big statue in the middle of the square.
\end{quote}

Built on \emph{jesm} 'body': a thing given a body.

\begin{cousinweb}
\emph{mojassame} was borrowed from Arabic \emph{mujassama} (EtymDB glosses it 'maquette'), a word built on \emph{jism} 'body'.
\end{cousinweb}

\begin{grammarlens}[title={What you've built}]
One more word for telling, explaining and describing at B1.
\end{grammarlens}

\subsection*{Your turn}
\begin{itemize}
\raggedright
  \item \emph{Say it:} \emph{mojassame}
  \item \emph{Say it:} \emph{mojassame}, once more
  \item \emph{Say it:} a sentence of your own with \emph{mojassame}
  \item \emph{Recall:} say the Persian for the outskirts, then the Persian for a tower, then say \emph{mojassame} again
\end{itemize}

\begin{tcolorbox}[breakable,colback=teal!4,colframe=teal!35!black,title={Before you move on}]
What does \fa{مجسمه} mean? (\textquotedblleft{}A statue\textquotedblright{}.) Say it once more.
\end{tcolorbox}
\section[\texorpdfstring{\fa{فواره} (favvâre)}{فواره (favvâre)}]{\fa{فواره} (favvâre) — a fountain}

\label{lesson:FA-C450-favvare}



\begin{quote}
Before the new one: say the Persian for a tower, then the Persian for a statue.
\end{quote}

\subsection*{You'll want to know: \fa{فواره}}
\textbf{\fa{فواره}} (\emph{favvâre}) — \textquotedblleft{}a fountain\textquotedblright{}.

\begin{quote}
\fa{بچه} \fa{ها} \fa{کنار} \fa{فواره} \fa{پارک} \fa{بازی} \fa{میکردند}.
\end{quote}

\begin{quote}
\emph{Bachche-hâ kenâr-e favvâre-ye pârk bâzi mikardand.}
\end{quote}

\begin{quote}
The children were playing by the fountain in the park.
\end{quote}

A spring in the ground is \emph{cheshme}; \emph{favvâre} sends water up.

\begin{cousinweb}
\emph{favvâre} was borrowed from Arabic \emph{fawwāra} (EtymDB), a word for water that gushes up.
\end{cousinweb}

\begin{grammarlens}[title={What you've built}]
One more word for telling, explaining and describing at B1.
\end{grammarlens}

\subsection*{Your turn}
\begin{itemize}
\raggedright
  \item \emph{Say it:} \emph{favvâre}
  \item \emph{Say it:} \emph{favvâre}, once more
  \item \emph{Say it:} a sentence of your own with \emph{favvâre}
  \item \emph{Recall:} say the Persian for a tower, then the Persian for a statue, then say \emph{favvâre} again
\end{itemize}

\begin{tcolorbox}[breakable,colback=teal!4,colframe=teal!35!black,title={Before you move on}]
What does \fa{فواره} mean? (\textquotedblleft{}A fountain\textquotedblright{}.) Say it once more.
\end{tcolorbox}
